% GENERATED FILE. Edit canonical lessons, then run npm run generate:books.
% canonical-source-hash: fnv1a64:2e1a24402171cd84
% canonical-lessons: FR-C376-faculte, FR-C376-inscrire, FR-C376-inscription, FR-C376-semestre, FR-C376-amphitheatre

\chapter{À l'université}
\label{ch:fr-a-l-universite}
\hlchaptermodality{376}

\begin{chapteropening}
\textbf{By the end of this chapter:} \emph{I can talk about starting university: the faculty, enrolling, enrolment, the semester and the lecture hall.}

Complete the last lesson of chapter 376: I can talk about starting university: the faculty, enrolling, enrolment, the semester and the lecture hall.
\end{chapteropening}

\section[la faculté]{la faculté — the faculty, the university; the ability}

\label{lesson:FR-C376-faculte}



\begin{quote}
Before the new one: say the French for the school report, then the French for the head teacher.
\end{quote}

\subsection*{You'll want to know: la faculté}
\textbf{la faculté} — \textquotedblleft{}the faculty, the university; the ability\textquotedblright{}.

\begin{quote}
Elle étudie à la faculté de médecine.
\end{quote}

\begin{quote}
\emph{She studies at the faculty of medicine.}
\end{quote}

In speech, \emph{la fac}: \emph{aller à la fac}. It can also mean an ability.

\begin{cousinweb}
\emph{faculté} is a learned borrowing from Latin \emph{facultās} \textquotedblleft{}ability, power\textquotedblright{} (Lewis \& Short), from \emph{facilis} \textquotedblleft{}easy, doable\textquotedblright{}.
\end{cousinweb}

\begin{grammarlens}[title={What you've built}]
One more word for telling, explaining and describing at B1.
\end{grammarlens}

\subsection*{Your turn}
\begin{itemize}
\raggedright
  \item \emph{Say it:} \emph{la faculté}
  \item \emph{Say it:} \emph{la faculté}, once more
  \item \emph{Say it:} a sentence of your own with \emph{la faculté}
  \item \emph{Recall:} say the French for the school report, then the French for the head teacher, then say \emph{la faculté} again
\end{itemize}

\begin{tcolorbox}[breakable,colback=teal!4,colframe=teal!35!black,title={Before you move on}]
What does la faculté mean? (\textquotedblleft{}The faculty, the university; the ability\textquotedblright{}.) Say it once more.
\end{tcolorbox}
\section[inscrire]{inscrire — to enrol, to register (mostly reflexive, s'inscrire); to write down}

\label{lesson:FR-C376-inscrire}



\begin{quote}
Before the new one: say the French for the head teacher, then the French for the faculty.
\end{quote}

\subsection*{You'll want to know: inscrire}
\textbf{inscrire} — \textquotedblleft{}to enrol, to register (mostly reflexive, s'inscrire); to write down\textquotedblright{}.

\begin{quote}
Je me suis inscrit à un cours d'anglais.
\end{quote}

\begin{quote}
\emph{I enrolled in an English course.}
\end{quote}

\emph{S'inscrire à} — to sign up for. Conjugated like \emph{écrire}: \emph{je m'inscris}.

\begin{cousinweb}
\emph{inscrire} is a learned borrowing from Latin \emph{īnscrībere} \textquotedblleft{}to write on, to inscribe\textquotedblright{} (EtymDB; Lewis \& Short), reshaped on \emph{écrire}; it is built on \emph{scrībere} \textquotedblleft{}to write\textquotedblright{}.
\end{cousinweb}

\begin{grammarlens}[title={What you've built}]
One more word for telling, explaining and describing at B1.
\end{grammarlens}

\subsection*{Your turn}
\begin{itemize}
\raggedright
  \item \emph{Say it:} \emph{inscrire}
  \item \emph{Say it:} \emph{inscrire}, once more
  \item \emph{Say it:} a sentence of your own with \emph{inscrire}
  \item \emph{Recall:} say the French for the head teacher, then the French for the faculty, then say \emph{inscrire} again
\end{itemize}

\begin{tcolorbox}[breakable,colback=teal!4,colframe=teal!35!black,title={Before you move on}]
What does inscrire mean? (\textquotedblleft{}To enrol, to register (mostly reflexive, s'inscrire); to write down\textquotedblright{}.) Say it once more.
\end{tcolorbox}
\section[l'inscription]{l'inscription — the registration, the enrolment}

\label{lesson:FR-C376-inscription}



\begin{quote}
Before the new one: say the French for the faculty, then the French for to enrol.
\end{quote}

\subsection*{You'll want to know: l'inscription}
\textbf{l'inscription} — \textquotedblleft{}the registration, the enrolment\textquotedblright{}.

\begin{quote}
Les inscriptions ferment le trente septembre.
\end{quote}

\begin{quote}
\emph{Registration closes on the thirtieth of September.}
\end{quote}

Also words carved in stone. Feminine.

\begin{cousinweb}
\emph{inscription} is a learned borrowing from Latin \emph{īnscrīptiō} \textquotedblleft{}a writing on, an inscription\textquotedblright{} (EtymDB; Lewis \& Short), from \emph{īnscrībere} \textquotedblleft{}to write on\textquotedblright{}.
\end{cousinweb}

\begin{grammarlens}[title={What you've built}]
One more word for telling, explaining and describing at B1.
\end{grammarlens}

\subsection*{Your turn}
\begin{itemize}
\raggedright
  \item \emph{Say it:} \emph{l'inscription}
  \item \emph{Say it:} \emph{l'inscription}, once more
  \item \emph{Say it:} a sentence of your own with \emph{l'inscription}
  \item \emph{Recall:} say the French for the faculty, then the French for to enrol, then say \emph{l'inscription} again
\end{itemize}

\begin{tcolorbox}[breakable,colback=teal!4,colframe=teal!35!black,title={Before you move on}]
What does l'inscription mean? (\textquotedblleft{}The registration, the enrolment\textquotedblright{}.) Say it once more.
\end{tcolorbox}
\section[le semestre]{le semestre — the semester, the half-year}

\label{lesson:FR-C376-semestre}



\begin{quote}
Before the new one: say the French for to enrol, then the French for the registration.
\end{quote}

\subsection*{You'll want to know: le semestre}
\textbf{le semestre} — \textquotedblleft{}the semester, the half-year\textquotedblright{}.

\begin{quote}
Le premier semestre commence en septembre.
\end{quote}

\begin{quote}
\emph{The first semester starts in September.}
\end{quote}

Six months; three months is \emph{un trimestre}.

\begin{cousinweb}
\emph{semestre} comes from Latin \emph{sēmēstris} \textquotedblleft{}of six months\textquotedblright{} (EtymDB), built from \emph{sex} \textquotedblleft{}six\textquotedblright{} and \emph{mēnsis} \textquotedblleft{}month\textquotedblright{} (Lewis \& Short).
\end{cousinweb}

\begin{grammarlens}[title={What you've built}]
One more word for telling, explaining and describing at B1.
\end{grammarlens}

\subsection*{Your turn}
\begin{itemize}
\raggedright
  \item \emph{Say it:} \emph{le semestre}
  \item \emph{Say it:} \emph{le semestre}, once more
  \item \emph{Say it:} a sentence of your own with \emph{le semestre}
  \item \emph{Recall:} say the French for to enrol, then the French for the registration, then say \emph{le semestre} again
\end{itemize}

\begin{tcolorbox}[breakable,colback=teal!4,colframe=teal!35!black,title={Before you move on}]
What does le semestre mean? (\textquotedblleft{}The semester, the half-year\textquotedblright{}.) Say it once more.
\end{tcolorbox}
\section[l'amphithéâtre]{l'amphithéâtre — the lecture hall}

\label{lesson:FR-C376-amphitheatre}



\begin{quote}
Before the new one: say the French for the registration, then the French for the semester.
\end{quote}

\subsection*{You'll want to know: l'amphithéâtre}
\textbf{l'amphithéâtre} — \textquotedblleft{}the lecture hall\textquotedblright{}.

\begin{quote}
Le cours a lieu dans le grand amphithéâtre.
\end{quote}

\begin{quote}
\emph{The lecture takes place in the big lecture hall.}
\end{quote}

Shortened to \emph{l'amphi}. Also an ancient arena. Masculine.

\begin{cousinweb}
\emph{amphithéâtre} comes from Latin \emph{amphitheātrum} (EtymDB), from Greek \emph{amphithéatron}, built from \emph{amphí} \textquotedblleft{}around\textquotedblright{} and \emph{théatron} \textquotedblleft{}a place for seeing\textquotedblright{}.
\end{cousinweb}

\begin{grammarlens}[title={What you've built}]
One more word for telling, explaining and describing at B1.
\end{grammarlens}

\subsection*{Your turn}
\begin{itemize}
\raggedright
  \item \emph{Say it:} \emph{l'amphithéâtre}
  \item \emph{Say it:} \emph{l'amphithéâtre}, once more
  \item \emph{Say it:} a sentence of your own with \emph{l'amphithéâtre}
  \item \emph{Recall:} say the French for the registration, then the French for the semester, then say \emph{l'amphithéâtre} again
\end{itemize}

\begin{tcolorbox}[breakable,colback=teal!4,colframe=teal!35!black,title={Before you move on}]
What does l'amphithéâtre mean? (\textquotedblleft{}The lecture hall\textquotedblright{}.) Say it once more.
\end{tcolorbox}
